% ======================================================================
% Abstract — one page, no section headings, max 3 references
% ======================================================================
\documentclass[12pt,a4paper]{article}

\usepackage[utf8]{inputenc}
\usepackage[T1]{fontenc}
\usepackage[english]{babel}
\usepackage{geometry}
\geometry{top=2.5cm, bottom=2.5cm, left=2.5cm, right=2.5cm}
\usepackage{amsmath}
\usepackage{amssymb}
\usepackage{mathptmx}        % Times New Roman text + math
\usepackage{setspace}
\usepackage{hyperref}
\usepackage{xcolor}
\usepackage{parskip}

% --- No indent, no extra space between paragraphs ---
\setlength{\parindent}{0pt}
\setlength{\parskip}{0pt}
\singlespacing

\definecolor{codeblue}{rgb}{0.1,0.2,0.6}
\hypersetup{colorlinks=true, linkcolor=codeblue, urlcolor=codeblue,
    pdftitle={Abstract: Inertia and Damping in Compass-Needle Lattices},
    pdfauthor={J. P. Sinnecker}}

% --- Dynamic 24 h compile-time stamp ---
\newcount\hours \newcount\minutes \newcount\temp
\hours=\time \divide\hours by 60
\minutes=\time \temp=\hours \multiply\temp by 60 \advance\minutes by -\temp
\def\compileTime{%
  \ifnum\hours<10 0\fi\the\hours:\ifnum\minutes<10 0\fi\the\minutes}

% --- Suppress page number ---
\pagestyle{empty}

% --- Compact title spacing ---
\usepackage{titling}
\setlength{\droptitle}{-2.8em}
\pretitle{\begin{center}\large\bfseries}
\posttitle{\par\end{center}\vspace{2pt}}
\preauthor{\begin{center}\normalsize\bfseries}
\postauthor{\par\end{center}\vspace{0pt}}
\predate{\begin{center}\small\itshape}
\postdate{\par\end{center}\vspace{6pt}}

\title{Influence of Inertia and Viscous Damping on the Magnetic Hysteresis
and First-Order Reversal Curve (FORC) Diagrams of Compass-Needle Lattices}
\author{J.~P.~Sinnecker\\[1pt]
\normalfont\small\itshape
Centro Brasileiro de Pesquisas F\'isicas (CBPF), Rio de Janeiro, Brazil}
\date{\today\ at \compileTime}

\begin{document}
\maketitle
\thispagestyle{empty}

% -----------------------------------------------------------------------
% BODY — single paragraph, no label, single-spaced, Times New Roman 12 pt
% -----------------------------------------------------------------------
\noindent
The collective behavior of two-dimensional arrays of magnetic dipoles has attracted
sustained interest in the context of artificial spin ice (ASI) and nanostructured
magnetic islands fabricated by nanolithography.
Classic thermodynamic studies of compass-needle lattices~\cite{prb1998} and the
introduction of First-Order Reversal Curve (FORC) diagrams~\cite{roberts2000}
established the conceptual tools now widely used to map coercive and interaction
field distributions in magnetically ordered arrays.
However, existing models rely on quasi-static or Monte Carlo approaches that omit
physical rotational inertia ($I$) and viscous damping ($b$), effectively assuming
an infinite-damping limit ($Q\!\to\!0$); variable FORC sweep rates further introduce
relaxation artifacts that distort coercive-field scaling and mask cooperative switching
cascades~\cite{ruta2017}.
We address these gaps with a Newtonian rotational molecular-dynamics model---integrated
by a Velocity-Verlet scheme---in which each needle obeys
$I\ddot{\theta}_{i} = \tau_{i}^{\text{dip}} + \tau_{i}^{\text{ext}} - b\dot{\theta}_{i}$,
with $I$ auto-calibrated from a rectangular blade plus a central cylindrical brass pivot.
Hysteresis loops and FORC diagrams are computed for square, triangular, and honeycomb
geometries across a wide range of quality factors $Q=\omega_{0}I/b$, and under a
strictly constant field sweep rate $R=|dB/dt|$ that eliminates rate-artifact
contamination.
Inertia delays switching, dynamically inflates the coercive field, and generates
phase-coherent oscillations at high $Q$; the resulting FORC broadening mimics
structural disorder.
Crucially, the mechanical damping ratio $b/I$ maps directly onto the phenomenological
Gilbert damping parameter $\alpha$ of thin-film nanomagnets, establishing this
macroscopic system as a transparent analog for nanolithography-fabricated magnetic
island arrays and providing quantitative guidance for interpreting FORC measurements
on nanopatterned samples.

% -----------------------------------------------------------------------
% REFERENCES — no heading, single-spaced, same font
% -----------------------------------------------------------------------
\vspace{4pt}
\begin{singlespace}
\begin{thebibliography}{9}
\setlength{\itemsep}{0pt}
\setlength{\parsep}{0pt}

\bibitem{prb1998}
E.~Olive and P.~Molho,
\emph{Thermodynamic study of a lattice of compass needles in dipolar interaction},
Phys.\ Rev.\ B \textbf{58}, 9238 (1998).

\bibitem{roberts2000}
A.~P.~Roberts, C.~R.~Pike, and K.~L.~Verosub,
\emph{First-order reversal curve diagrams: A new tool for characterizing the
magnetic properties of natural samples},
J.\ Geophys.\ Res.\ Solid Earth \textbf{105}, 28461 (2000).

\bibitem{ruta2017}
S.~Ruta \emph{et al.},
\emph{First order reversal curves and intrinsic parameter determination
for magnetic materials},
Sci.\ Rep.\ \textbf{7}, 45218 (2017).

\end{thebibliography}
\end{singlespace}

\end{document}
